% se não for usar a quarta palavra chave, deixar o campo vazio: {}
\palavraschaves{Covid-19}
{Modelagem baseada em agentes}
{Programação paralela em GPU}
{CUDA}



\pretextualchapter{Resumo}

\noindent
As simulações epidemiológicas baseadas em agentes permitem representar a evolução de uma doença em uma população levando em conta as características individuais e espaciais dos indivíduos, porém, para áreas com grande número de indivíduos e diversos parâmetros, demandam elevado poder computacional. Este trabalho apresenta a paralelização, utilizando CUDA, de um modelo baseado em agentes do tipo SEIR estendido para a Covid-19, aplicado a quatro áreas urbanas brasileiras: São Paulo, Rocinha, Brasília e Manaus. A implementação em GPU é validada contra a implementação serial de referência sob os mesmos parâmetros, e o seu desempenho é avaliado por meio da aceleração e da eficiência. Os resultados mostram que a implementação paralela reproduz as curvas epidêmicas da implementação serial, com erro quadrático médio inferior a 1,2 ponto percentual, valor compatível com a flutuação estatística do método de Monte Carlo, e que a aceleração obtida cresce com o tamanho da rede, atingindo cerca de 56 vezes para a maior cidade. A análise de desempenho indica que a simulação é limitada pela largura de banda de memória, e uma otimização da estrutura de dados que mantém os acessos na memória cache reduz significativamente o tempo de execução.

\imprimirchaves % linha em branco antes
